\section{Preliminaries}
\label{sec:overview}
% ============================================================

\subsection{Invertible Bloom Lookup Tables}

Invertible Bloom Lookup Tables (IBLTs) were introduced by Goodrich and
Mitzenmacher~\cite{GM11}. An IBLT is a variant of Bloom filter that supports,
in addition to inserting and deleting elements, the ability to \emph{list} all
stored elements. Unlike a standard Bloom filter---which only answers membership
queries with a configurable false-positive rate and requires $\approx 60\times$
expansion for negligible error---an IBLT stores key-value pairs in its cells
and achieves listing with expansion factor $1.5$--$4.5\times$ (for $k = 4$ hash
functions). This makes IBLTs uniquely suited for set reconciliation and private
set union, where the goal is to recover the entire set rather than answer
individual membership queries.

An IBLT is parameterized by $\prm = (M, k, \ell, H)$: a modulus $M \in
\N_{\geq 2}$ (typically $2^{64}$), $k$ hash functions $H = \{h_i : \U \to
[\ell]\}_{i \in [k]}$ each mapping to a distinct subtable of size $\ell$, and
two $k \times \ell$ matrices $\mathsf{sum}$ (values in $\Z_M$) and
$\mathsf{cnt}$ (non-negative integers). We write $m = k\ell$ for the total
number of cells.

The three core operations are given in Algorithms~\ref{alg:encode}--\ref{alg:list}.

\begin{algorithm}[ht]
\caption{$\Encode_{\prm}(X)$ --- IBLT encoding.}
\label{alg:encode}
\begin{algorithmic}
\STATE \textbf{Input:} $\prm = (M, k, \ell, H)$, set $X \subseteq \U$, empty IBLT $(\mathsf{sum}, \mathsf{cnt})$.
\STATE \textbf{Output:} $\mathsf{IBLT} = (\mathsf{sum}, \mathsf{cnt}, \prm)$ containing all $x \in X$.
\STATE Initialize $\mathsf{sum}[i][j] \leftarrow 0$, $\mathsf{cnt}[i][j] \leftarrow 0$ for all $i \in [k], j \in [\ell]$.
\FOR{each $x \in X$}
  \FOR{$i \in [k]$}
    \STATE $j \leftarrow h_i(x)$
    \STATE $\mathsf{sum}[i][j] \leftarrow \mathsf{sum}[i][j] + x \pmod M$
    \STATE $\mathsf{cnt}[i][j] \leftarrow \mathsf{cnt}[i][j] + 1$
  \ENDFOR
\ENDFOR
\end{algorithmic}
\end{algorithm}

\begin{algorithm}[ht]
\caption{$\DeleteSet_{\prm}(\mathsf{IBLT}, X)$ --- batch deletion.}
\label{alg:deleteset}
\begin{algorithmic}
\STATE \textbf{Input:} $\mathsf{IBLT} = (\mathsf{sum}, \mathsf{cnt}, \prm)$ and set $X \subseteq \U$.
\STATE \textbf{Output:} Updated $\mathsf{IBLT} = (\mathsf{sum}, \mathsf{cnt}, \prm)$ with all $x \in X$ removed.
\FOR{each $x \in X$}
  \FOR{$i \in [k]$}
    \STATE $j \leftarrow h_i(x)$
    \STATE $\mathsf{sum}[i][j] \leftarrow \mathsf{sum}[i][j] - x \pmod M$
    \STATE $\mathsf{cnt}[i][j] \leftarrow \mathsf{cnt}[i][j] - 1$
  \ENDFOR
\ENDFOR
\end{algorithmic}
\end{algorithm}

IBLT encoding is \emph{linear}: $\Encode(A \cup B) = \Encode(A) \oplus
\Encode(B)$ and $\Encode(A \setminus B) = \Encode(A) \ominus \Encode(B)$,
where $\oplus$/$\ominus$ denote cell-wise modular addition/subtraction.

A cell with $\mathsf{cnt}[i][j] = 1$ is a \emph{singleton}: its sum equals
the unique element stored there.

\begin{algorithm}[ht]
\caption{$\List(\mathsf{IBLT})$ --- full listing (peeling).}
\label{alg:list}
\begin{algorithmic}
\STATE \textbf{Input:} $\mathsf{IBLT} = (\mathsf{sum}, \mathsf{cnt}, \prm)$.
\STATE \textbf{Output:} Set $V$ of all encoded elements, or $\bot$ on failure.
\STATE $V \leftarrow \emptyset$
\WHILE{there exists $(i,j)$ with $\mathsf{cnt}[i][j] = 1$}
  \STATE $x \leftarrow \mathsf{sum}[i][j]$
  \STATE $V \leftarrow V \cup \{x\}$
  \STATE $\DeleteSet_{\prm}(\mathsf{IBLT}, \{x\})$
\ENDWHILE
\IF{$\mathsf{cnt}[i][j] = 0$ for all $(i,j)$}
  \STATE \textbf{output} $V$
\ELSE
  \STATE \textbf{output} $\bot$
\ENDIF
\end{algorithmic}
\end{algorithm}

\begin{theorem}[IBLT Listing~\cite{GM11,PT26}]\label{thm:iblt}
For $m = k\ell > (c_k + \varepsilon)\tau$ with $n \leq \tau$ elements,
$\Pr[\List\text{ fails}] = O(\tau^{-k+2})$. With $k = 4$, expansion
$m/\tau \in [1.5, 4.5]$, failure $\leq 2^{-\lambda}$.
\end{theorem}

In the above, $c_k$ is the threshold constant for $2$-core emergence ($c_3
\approx 1.222$, $c_4 \approx 1.295$, $c_5 \approx 1.425$~\cite{PT26}). If $\List$ terminates with non-singleton
cells remaining (a $2$-core), listing has failed. The expansion factor $m/n$
must exceed $c_k$ to guarantee an empty $2$-core with overwhelming probability.

\subsection{Union Peelability}

Piske and Trieu~\cite{PT26} observed that the union of two IBLTs can be peeled
without materializing the combined table. For $\mathsf{IBLT}_0 =
\Encode(X_0)$ and $\mathsf{IBLT}_1 = \Encode(X_1)$, the virtual union is
$\mathsf{IBLT}_\cup = \mathsf{IBLT}_0 \oplus \mathsf{IBLT}_1$.

\begin{definition}[$\uPeel$~\cite{PT26}]\label{def:upeel}
\footnotesize
\[
\uPeel(\mathsf{IBLT}_0, \mathsf{IBLT}_1, i, j) = \begin{cases}
\mathsf{sum}_0[i][j], & \mathsf{cnt}_0{=}1 \land \mathsf{cnt}_1{=}0 \\
\mathsf{sum}_1[i][j], & \mathsf{cnt}_0{=}0 \land \mathsf{cnt}_1{=}1 \\
\mathsf{sum}_0[i][j], & \mathsf{cnt}_0{=}\mathsf{cnt}_1{=}1
                       \land \mathsf{sum}_0{=}\mathsf{sum}_1 \\
\bot, & \text{otherwise}
\end{cases}
\]
\end{definition}

Intuitively, $\uPeel$ distinguishes three cases: the peeled element is unique
to $P_0$ (case~1), unique to $P_1$ (case~2), or shared with equal sums (case~3).
Secure evaluation uses OT for sum retrieval and bOPRF for equality comparison.

\begin{theorem}[Union Peel Equivalence~\cite{PT26}]\label{thm:upeel}
For union-peelable IBLTs: $\Peel(\mathsf{IBLT}_\cup, i, j) =
\uPeel(\mathsf{IBLT}_0, \mathsf{IBLT}_1, i, j)$.
\end{theorem}

\begin{center}
\fbox{\footnotesize\begin{minipage}{0.96\columnwidth}
\sloppy
\smallskip
\noindent\textbf{$\Pi_{\mathsf{UnionPeel}}$ --- secure union peeling~\cite{PT26}}
\smallskip

\noindent On input $Q \subseteq [k] \times [\ell]$ (cells to scan) and IBLTs
$(\mathsf{IBLT}_C, \mathsf{IBLT}_S)$:
\begin{enumerate}
\item For each $(i,j) \in Q$: evaluate $\uPeel(\mathsf{IBLT}_C, \mathsf{IBLT}_S, i, j)$
      via $\cF_{\mathsf{OT}}$ and $\cF_{\mathsf{OPRF}}$.
\item Collect recovered elements $V \leftarrow \{v \mid \uPeel(\cdot) \neq \bot\}$.
\item Output $V$ to both parties.
\end{enumerate}
$\Pi_{\mathsf{UnionPeel}}$ UC-realizes $\cF_{\mathsf{UnionPeel}}$ in the
$(\cF_{\mathsf{OPRF}}, \cF_{\mathsf{OT}})$-hybrid model against
semi-honest adversaries.
\smallskip
\end{minipage}}  % closes minipage, fbox, and setlength group
\captionof{figure}{Sub-protocols used in F-PSU~\cite{PT26}.}
\label{fig:unionpeel}
\end{center}
\begin{theorem}[Union Peel Security~\cite{PT26}]\label{thm:unionpeel-sec}
$\Pi_{\mathsf{UnionPeel}}$ UC-realizes $\cF_{\mathsf{UnionPeel}}$
(Fig.~\ref{fig:unionpeel}) in the $(\cF_{\mathsf{OPRF}},
\cF_{\mathsf{OT}})$-hybrid model against semi-honest adversaries, achieving
static PSU on $2^{20}$ elements in under 3~seconds.
\end{theorem}

\subsection{Cryptographic Primitives}

\textbf{Symmetric encryption.} We use a symmetric encryption scheme
$\mathsf{SE} = (\mathsf{SE.Enc}, \mathsf{SE.Dec})$ with 256-bit keys.
Our implementation instantiates $\mathsf{SE}$ with ChaCha20~\cite{Ber08}
(256-bit key, 96-bit nonce, $C = M \oplus \mathsf{ChaCha20.keystream}(K, N)$).
The IBLT ciphertexts are encrypted under per-epoch keys derived from the
PRF chain.

\textbf{HMAC-SHA256~\cite{BCK96}.} We use HMAC-SHA256 as a secure PRF:
$F: \{0,1\}^\kappa \times \{0,1\}^* \to \{0,1\}^\kappa$. Key evolution uses a
fixed domain separator: $K_t \leftarrow F(K_{t-1}, \text{``evolve''})$.

\begin{lemma}[PRF Key One-Wayness]\label{lem:prf-ow}
If $F$ is a secure PRF, then for any PPT $\cA$ and fixed $c$:
$\Pr[\cA(F(K, c)) = K : K \sample \{0,1\}^\kappa] \leq \negl(\kappa)$.
\end{lemma}

\textbf{VOLE and OT extension.} Vector Oblivious Linear Evaluation (VOLE
enables two parties to obtain additive shares of $\Delta \cdot \vec{x}$,
realizable from $\kappa$ base OTs with sublinear communication~\cite{RS21,RR22};
silent VOLE~\cite{RRT23} eliminates the base OTs entirely. The VOLE state is
reusable across all epochs, amortizing the one-time setup cost.

\textbf{Batched OPRF.} The KKRT16/RS21/RR22 framework~\cite{KKRT16,RS21,RR22,HKL+24}
constructs batched OPRF from VOLE and OKVS: the receiver encodes queries into
an OKVS $P$; VOLE shares $\Delta \cdot P$; OPRF outputs are defined as
$F_\Delta(x) = H(\Delta \cdot \mathsf{Decode}(P, x))$.

The ideal 1-out-of-2 OT functionality $\cF_{\mathsf{OT}}$ is defined in
Fig.~\ref{fig:f-ot}.

\begin{center}
\fbox{\begin{minipage}{0.95\columnwidth}
\smallskip
\noindent\textbf{Functionality $\cF_{\mathsf{OT}}$}

\noindent\textbf{Parameters:} Sender $S$ holds two messages
$(m_0, m_1) \in \{0,1\}^\kappa$; receiver $R$ holds a choice bit
$b \in \{0,1\}$.

\smallskip\noindent\textbf{Behavior:}
\begin{enumerate}
\item Receive $(m_0, m_1)$ from $S$ and $b$ from $R$. If either party
sends $\bot$, output $\bot$ to both and halt.
\item Send $m_b$ to $R$ and nothing to $S$.
\end{enumerate}
\smallskip
\end{minipage}}
\captionof{figure}{The ideal 1-out-of-2 Oblivious Transfer functionality.}
\label{fig:f-ot}
\end{center}

The ideal OPRF functionality $\cF_{\mathsf{OPRF}}$ is defined in
Fig.~\ref{fig:f-oprf}.

\begin{center}
\fbox{\begin{minipage}{0.95\columnwidth}
\smallskip
\noindent\textbf{Functionality $\cF_{\mathsf{OPRF}}$}

\noindent\textbf{Parameters:} Sender $S$ holds PRF key $k \in \{0,1\}^\kappa$;
receiver $R$ holds input $x \in \U$.

\smallskip\noindent\textbf{Behavior:}
\begin{enumerate}
\item Receive $k$ from $S$ and $x$ from $R$. If either party sends $\bot$,
output $\bot$ to both and halt.
\item Compute $y \leftarrow F(k, x)$. Send $y$ to $R$ and nothing to $S$.
\end{enumerate}
\smallskip
\end{minipage}}
\captionof{figure}{The ideal Oblivious PRF functionality.}
\label{fig:f-oprf}
\end{center}

The standard (static) PSU functionality $\cF_{\mathsf{PSU}}$ is defined
in Fig.~\ref{fig:f-psu}.

\begin{center}
\fbox{\begin{minipage}{0.95\columnwidth}
\smallskip
\noindent\textbf{Functionality $\cF_{\mathsf{PSU}}$}

\noindent\textbf{Parameters:} Two parties $P_0, P_1$ with element domain $\U$.

\smallskip\noindent\textbf{Behavior:}
\begin{enumerate}
\item Receive set $X$ from $P_0$ and $Y$ from $P_1$. If either is $\bot$,
output $\bot$ to both and halt.
\item Send $X \cup Y$ to both parties.
\end{enumerate}
\smallskip
\end{minipage}}
\captionof{figure}{The ideal (static) Private Set Union functionality.}
\label{fig:f-psu}
\end{center}

\subsection{UC Framework}

We use the Universal Composability framework~\cite{Can01} with adaptive
corruptions. A protocol $\Pi$ UC-realizes an ideal functionality $\cF$ if for
every PPT adversary $\cA$ there exists a PPT simulator $\cS$ such that no PPT
environment $\cZ$ can distinguish the real execution from the ideal execution.
We operate in the $(\cF_{\mathsf{OPRF}}, \cF_{\mathsf{OT}})$-hybrid model,
where the protocol invokes the ideal functionalities defined in
Figs.~\ref{fig:f-ot},~\ref{fig:f-oprf}. The standard (static) PSU
functionality $\cF_{\mathsf{PSU}}$ (Fig.~\ref{fig:f-psu}) and the
union-peeling functionality $\cF_{\mathsf{UnionPeel}}$
(Fig.~\ref{fig:unionpeel}) provide the building blocks. The ideal
functionality for our protocol, $\cF_{\mathsf{fpsu}}$, is specified in
Section~\ref{sec:protocol} (Fig.~\ref{fig:f-fpsu}).

